\section{Named-Competitor Comparison Scaffold}
\label{sec:fhe-named-competitors}
%% ============================================================================

To make the Lux-FHE claims falsifiable we publish here a comparison-table
scaffold against the most widely benchmarked open-source FHE libraries
that share an open, reproducible benchmark surface.
\textbf{Lux columns are \texttt{[TBD-measure]}} until the same circuits run
on the same hardware as the competitor's published benchmark.
\textbf{Competitor numbers are taken verbatim from the cited
documentation / paper}; the source for every cell is listed in the
last column.

\subsection{Competitors}

\begin{itemize}[nosep]
  \item \textbf{TFHE-CGGI reference}~\cite{cggi} --- the original
        Chillotti-Gama-Georgieva-Iza\-bach\`ene construction; the
        academic baseline for binary-gate / programmable-bootstrapping FHE.
  \item \textbf{OpenFHE}~\cite{albadawi2022openfhe} --- C++ library
        carrying BFV, BGV, CKKS, FHEW/TFHE; community standard for
        academic comparisons.
  \item \textbf{VeloFHE / OpenFHE-GPU}~\cite{velofhe2024} --- GPU-accelerated
        OpenFHE fork; the open GPU baseline.
  \item \textbf{Microsoft SEAL}~\cite{microsoftseal} --- BFV/CKKS;
        canonical pure-CPU baseline for arithmetic FHE.
\end{itemize}

\subsection{Boolean-Gate Bootstrapping Latency}

\begin{table}[h]
\centering
\small
\begin{tabular}{@{}llrrl@{}}
\toprule
\textbf{Library} & \textbf{Hardware} & \textbf{Gate latency} & \textbf{PBS / s} & \textbf{Source} \\
\midrule
\textbf{Lux-FHE (CPU)} & EPYC          & \texttt{[TBD-measure]} & \texttt{[TBD-measure]} & this paper \\
\textbf{Lux-FHE (CPU)} & Apple M2      & \texttt{[TBD-measure]} & \texttt{[TBD-measure]} & this paper \\
\textbf{Lux-FHE (Metal)} & M1/M2/M3 GPU & \texttt{[TBD-measure]} & \texttt{[TBD-measure]} & this paper, \cite{luxevmgpu} \\
\textbf{Lux-FHE (CUDA)} & 1$\times$H100 & \texttt{[TBD-measure]} & \texttt{[TBD-measure]} & this paper \\
\midrule
TFHE-CGGI ref. (CPU)      & i9-9900K               & $\sim$13\,ms         & $\sim$77 / core           & \cite{cggi} \\
OpenFHE FHEW              & x86\_64 (48-thread)    & $\sim$200\,ms      & $\sim$240 / 48 cores       & \cite{velofhe2024} \\
OpenFHE FHEW (GPU, VeloFHE) & RTX 4090            & $\sim$0.5\,ms     & $\sim$2{,}000             & \cite{velofhe2024} \\
\bottomrule
\end{tabular}
\caption{Boolean-gate (programmable bootstrapping) latency for FHE
backends. Numbers in the bottom block are quoted verbatim from the
cited source; Lux rows remain \texttt{[TBD-measure]}.}
\label{tab:fhe-pbs-comparison}
\end{table}

\subsection{Library Coverage Matrix}

\begin{table}[h]
\centering
\small
\begin{tabular}{@{}lccccl@{}}
\toprule
\textbf{Library} & \textbf{Binary FHE} & \textbf{Arith. FHE} & \textbf{GPU} & \textbf{Threshold} & \textbf{Source} \\
\midrule
Lux-FHE        & yes       & CKKS / BFV (Lattigo fork) & Metal / CUDA       & yes (LSS-Pulsar) & \cite{lattigov7} \\
OpenFHE        & yes       & BFV/BGV/CKKS      & via VeloFHE fork   & no                & \cite{albadawi2022openfhe,velofhe2024} \\
Microsoft SEAL & no        & BFV/CKKS          & no (CPU-only)      & no                & \cite{microsoftseal} \\
\bottomrule
\end{tabular}
\caption{Feature-coverage matrix for FHE backends in this scaffold.}
\label{tab:fhe-coverage}
\end{table}

\subsection{Reproducibility (what needs to be measured)}
\label{sec:fhe-repro}

\begin{itemize}[nosep]
  \item \textbf{Hardware.} Three rows, each repeated on each backend:
        (1) AWS \texttt{hpc7a.96xlarge} (192-core EPYC);
        (2) Apple M2 8-perf / 4-eff cores;
        (3) single H100 plus single Apple M3 Max (for the Lux Metal path).
  \item \textbf{Software pin.} \texttt{OpenFHE} v1.1.4 (with VeloFHE
        backend at HEAD); \texttt{SEAL} 4.1.1; Lux-FHE at the same
        git SHA as this paper's freeze tag.
  \item \textbf{Circuits.} Four operations on which cross-library
        reference numbers exist in the public literature today:
        (1) 4-bit boolean AND with PBS,
        (2) 8-bit integer multiply, (3) 32-bit equality test,
        (4) CKKS slot rotation.
  \item \textbf{Commands.}
  \begin{verbatim}
    # Lux-FHE
    go test -run=. -bench=BenchmarkBootstrap -benchtime=10s ./fhe/...
    # OpenFHE
    ./build/benchmark/binfhe-benchmark
    # SEAL
    ./build/bin/sealbench
  \end{verbatim}
  \item \textbf{Decision rule.} A Lux-FHE win against an open
        baseline requires that the \emph{same parameter set}
        (TUniform~4-bit, q=2$^{64}$) yields lower wall-clock PBS
        latency on the same hardware, in both single-core and
        full-machine configurations. Anything else is a not-equal
        comparison.
\end{itemize}

%% ============================================================================
